\section{Background}\label{sec:teoria}

\subsection{Transfer Learning}
The first layers of a convolutional network learn edges and textures that are useful for almost any visual task, while the last layers specialize in the source classes~\cite{yosinski}. It is therefore possible to reuse a network trained on ImageNet~\cite{imagenet}, replace its head, and train only part of it. Freezing more layers reduces overfitting when data are scarce, and unfreezing them allows the features to adapt to the new domain.

ResNet-50~\cite{resnet} learns residuals $\mathcal{F}(x)$ and outputs $\mathcal{F}(x)+x$, which makes it possible to train deep networks without gradient degradation. MobileNetV3~\cite{mobilenet} uses depthwise separable convolutions, inverted residual blocks, and channel attention to reduce operations and parameters.

\subsection{Connectionist Temporal Classification}
Let $y^t_k$ be the probability of symbol $k$ at step $t$, over an alphabet extended with the blank $\epsilon$. A path $\pi$ assigns one symbol to each step, and the function $\mathcal{B}$ reduces it to text by merging repetitions and removing blanks. CTC~\cite{ctc} defines the probability of a transcription $\ell$ as the sum over all of its paths:
\begin{equation}
p(\ell\mid x)=\sum_{\pi:\,\mathcal{B}(\pi)=\ell}\ \prod_{t=1}^{T} y^t_{\pi_t}.
\end{equation}
The sum is computed by dynamic programming in $O(T|\ell|)$, and the loss is $-\ln p(\ell\mid x)$. Training thus requires no labels for when each letter is spoken. The blank makes it possible to represent double letters: producing ``ll'' requires at least one $\epsilon$ between them.

\subsection{Decoding}
Greedy decoding takes the most probable symbol at each step. It is fast, but it ignores that several paths add up to the same text. Prefix beam search~\cite{beam} sums those paths and admits a language model, with the criterion
\begin{equation}
\hat{\ell}=\arg\max_{\ell}\ \ln p(\ell\mid x)+\alpha\ln p_{LM}(\ell)+\beta|\ell|,
\end{equation}
where $\alpha$ weights the language model and $\beta$ offsets its preference for short texts.

\subsection{DeepSpeech2}
DeepSpeech2~\cite{ds2} maps a spectrogram to character probabilities using 2D convolutions, which capture local time--frequency patterns and reduce resolution, followed by bidirectional recurrent layers, which provide context in both directions. It is the architecture studied in the course~\cite{keras-ctc}.

\subsection{Metrics}
Classification uses accuracy, precision, recall, and F1. Macro F1 averages the classes without weighting, so a minority class counts as much as the others. Speech uses the edit distance between reference and hypothesis:
\begin{equation}
\mathrm{WER}=\frac{S+D+I}{N},
\end{equation}
where $S$, $D$, and $I$ are word substitutions, deletions, and insertions, and $N$ is the number of reference words. CER is the equivalent over characters. WER penalizes an entire word for a single wrong letter, so it is usually considerably higher.
